\chapter{Théorie des catégories}

\section{Catégories}

\begin{definition}
Une \emph{catégorie} $\mathscr{C}$ consiste en :

\begin{itemize}
\item une collection $\text{ob}(\mathscr{C})$ d'\textbf{objets} ;
\item pour chaque $A, B \in \text{ob}(\mathscr{C})$, une collection $\mathscr{C}(A, B)$ de \textbf{morphismes} (ou \textbf{flèches} ou \textbf{applications}) de $A$ vers $B$ ;
\item pour chaque $A, B, C \in \text{ob}(\mathscr{C})$, une fonction
$$\mathscr{C}(B, C) \times \mathscr{C}(A, B) \longrightarrow \mathscr{C}(A, C)$$
$$(g, f) \longmapsto g \circ f,$$
appelée \textbf{composition} ;
\item pour chaque $A \in \text{ob}(\mathscr{C})$, un élément $1_A$ de $\mathscr{C}(A, A)$, appelé \textbf{identité} sur $A$,
\end{itemize}

satisfaisant les axiomes suivants :

\begin{itemize}
\item \textbf{associativité} : pour chaque $f \in \mathscr{C}(A, B)$, $g \in \mathscr{C}(B, C)$ et $h \in \mathscr{C}(C, D)$, nous avons $(h \circ g) \circ f = h \circ (g \circ f)$ ;
\item \textbf{lois d'identité} : pour chaque $f \in \mathscr{C}(A, B)$, nous avons $f \circ 1_A = f = 1_B \circ f$.
\end{itemize}
\end{definition}

\section{Foncteurs}
